\section{Quản lý dự án và phân công nhiệm vụ}

Để đảm bảo quá trình nghiên cứu và phát triển hệ thống diễn ra xuyên suốt, đúng tiến độ và đạt được các yêu cầu kỹ thuật đã đề ra, việc thiết lập một quy trình quản lý dự án bài bản là yếu tố tiên quyết. Chương này trình bày chi tiết về phương pháp luận quản lý được áp dụng, hệ sinh thái công cụ hỗ trợ, cách thức tổ chức nhóm và lộ trình thực hiện đồ án.

\subsection{Phương pháp quản lý dự án}

Với đặc thù của một đồ án phần mềm có kiến trúc phức tạp, bao gồm nhiều phân hệ tương tác lẫn nhau (Web Dashboard phân tích dữ liệu giao thông, Ứng dụng tra cứu cho người dân, Backend Server và Data Warehouse), nhóm quyết định áp dụng phương pháp phát triển phần mềm linh hoạt \textbf{Agile}, kết hợp với khung làm việc \textbf{Scrum}.

Việc áp dụng Agile/Scrum mang lại những lợi thế cụ thể cho nhóm:

\begin{itemize}
    \item \textbf{Tính thích ứng cao:} Cho phép nhóm dễ dàng điều chỉnh các tính năng (ví dụ: tinh chỉnh biểu đồ phân tích BI, thay đổi logic tính toán chỉ số TTI, BI) dựa trên kết quả nghiệm thu từng phần mà không làm đứt gãy toàn bộ cấu trúc dự án.
    \item \textbf{Phát triển lặp (Iterative Development):} Quá trình phát triển được chia nhỏ thành các chu kỳ ngắn gọi là Sprint (kéo dài từ 1 đến 2 tuần). Cuối mỗi Sprint, nhóm sẽ cho ra mắt một phần mềm có khả năng chạy được (shippable product) để đánh giá và kiểm thử ngay lập tức.
    \item \textbf{Tối ưu hóa giao tiếp:} Tổ chức các buổi họp ngắn (Sync-up meeting) từ 2-3 lần/tuần để cập nhật tiến độ, gỡ rối các rào cản kỹ thuật (bottlenecks) giữa các thành viên.
\end{itemize}

\textbf{Quy trình triển khai Scrum thực tế của nhóm:}
\begin{enumerate}
    \item \textbf{Product Backlog:} Tập hợp toàn bộ các tính năng cần có của hệ thống (ví dụ: Chức năng xem bản đồ nhiệt, Tính năng gợi ý lộ trình, Module ETL xử lý dữ liệu lịch sử).
    \item \textbf{Sprint Planning:} Lựa chọn các task từ Backlog đưa vào Sprint hiện tại dựa trên mức độ ưu tiên.
    \item \textbf{Sprint Execution:} Hiện thực hóa các tính năng bằng mã nguồn.
    \item \textbf{Sprint Review \& Retrospective:} Tổng kết, ghép nối (integrate) các module, kiểm tra lỗi (bug) và rút kinh nghiệm cho chu kỳ tiếp theo.
\end{enumerate}

\subsection{Công cụ quản lý và Môi trường làm việc nhóm}

Để hiện thực hóa phương pháp Agile, nhóm đã lựa chọn và đồng bộ hóa một hệ sinh thái công cụ quản lý chuyên nghiệp, bám sát với tiêu chuẩn của các doanh nghiệp phần mềm hiện nay:

\begin{itemize}
    \item \textbf{Quản lý mã nguồn (Source Code Management):} Sử dụng \textbf{Git} và nền tảng \textbf{GitHub} để lưu trữ mã nguồn. Áp dụng chiến lược Git Flow cơ bản, chia tách các nhánh \texttt{main} (môi trường production), \texttt{develop} (môi trường tích hợp) và các nhánh \texttt{feature/*} cho từng tính năng riêng biệt để tránh xung đột mã nguồn.
    \item \textbf{Quản lý công việc (Task Management):} Khởi tạo không gian làm việc trên \textbf{Linear} theo mô hình bảng Kanban. Các công việc được thiết lập dưới dạng thẻ (card) và di chuyển qua các cột trạng thái: \textit{To Do}, \textit{In Progress}, \textit{In Review}, và \textit{Done}.
    \item \textbf{Môi trường giao tiếp:} Sử dụng \textbf{Discord} làm kênh trao đổi thông tin chính, chia các channel (kênh) riêng biệt cho từng luồng công việc như \texttt{\#frontend-ui}, \texttt{\#backend-api}, và \texttt{\#deployment-alerts}.
    \item \textbf{Môi trường phát triển và Tích hợp:} Chuẩn hóa môi trường cục bộ (local environment) giữa các thành viên bằng \textbf{Docker}, đảm bảo tính nhất quán từ môi trường phát triển (Development) đến môi trường triển khai (Production). Sử dụng nền tảng chia sẻ tài liệu API (như Postman/Swagger) để Frontend và Backend có thể làm việc độc lập.
\end{itemize}

\subsection{Phân công vai trò và chi tiết nhiệm vụ}

Dựa trên thế mạnh cốt lõi và định hướng chuyên môn của từng thành viên, khối lượng công việc của đồ án được phân bổ một cách chuyên môn hóa nhưng vẫn đảm bảo tính phối hợp chặt chẽ (Cross-functional). Cụ thể như sau:

\vspace{0.5cm}
\noindent\textbf{Sinh viên 1: Hồ Trần Ngọc Liêm (MSSV: 2211835)}
\begin{itemize}
    \item \textbf{Vai trò:} Software Fullstack Developer.
    \item \textbf{Nhiệm vụ chi tiết:}
    \begin{itemize}
        \item Chịu trách nhiệm xây dựng kiến trúc giao diện, tối ưu hóa Trải nghiệm người dùng (UX) và Giao diện người dùng (UI) cho Web Dashboard điều hành và App người dân.
        \item Xây dựng và phát triển Frontend sử dụng thư viện React kết hợp với Ant Design để tạo ra các biểu đồ phân tích trực quan (Recharts), bản đồ số (Mapbox) xử lý khối lượng dữ liệu lớn.
        \item Đảm nhiệm vai trò theo dõi tiến độ chung và phối hợp triển khai hệ thống (Deployment).
    \end{itemize}
\end{itemize}

\vspace{0.5cm}
\noindent\textbf{Sinh viên 2: Đinh Nguyễn Việt Khiêm (MSSV: 2211565)}
\begin{itemize}
    \item \textbf{Vai trò:} Data Engineer (Hạ tầng dữ liệu).
    \item \textbf{Nhiệm vụ chi tiết:}
    \begin{itemize}
        \item Nghiên cứu và chỉnh sửa, tối ưu hoá cấu trúc Data Warehouse từ Giai đoạn 1 để lưu trữ luồng dữ liệu giao thông khổng lồ, phù hợp với Giai đoạn 2.
        \item Xử lý và làm sạch dữ liệu trạng thái giao thông, chuẩn hóa dữ liệu từ TomTom và các nguồn khác để đảm bảo tính nhất quán và chính xác. Đánh giá chất lượng dữ liệu đầu vào, kiểm tra lỗi, tìm kiếm nguồn dữ liệu thay thế khi cần thiết.
        \item Xây dựng và tối ưu hóa luồng ETL trên hạ tầng Azure Virtual Machine.
    \end{itemize}
\end{itemize}

\vspace{0.5cm}
\noindent\textbf{Sinh viên 3: Nguyễn Thành Dũng (MSSV: 2210589)}
\begin{itemize}
    \item \textbf{Vai trò:} AI/Algorithm Developer.
    \item \textbf{Nhiệm vụ chi tiết:}
    \begin{itemize}
        \item Nghiên cứu, xây dựng và tích hợp mô hình phân tích, dự báo điểm nghẽn giao thông.
        \item Chịu trách nhiệm xử lý các bài toán nghiệp vụ cốt lõi (tính toán các chỉ số TTI, BI, PTI theo chuẩn quốc tế).
        \item Thiết kế kịch bản kiểm thử (Load test, API Latency) và thực hiện đánh giá hiệu năng hệ thống toàn diện để đảm bảo tính ổn định của mô hình.
    \end{itemize}
\end{itemize}

\subsection{Tiến độ thực hiện đồ án}

Quá trình thực hiện đồ án được chia thành các giai đoạn chính (Milestones), trải dài trong suốt 15 tuần của học kỳ. Bảng \ref{tab:tiendo} dưới đây mô tả chi tiết các đầu mục công việc, kết quả kỳ vọng và người phụ trách cho từng giai đoạn.

\begin{table}[H]
    \centering
    \caption{Chi tiết các mốc thời gian thực hiện đồ án}
    \label{tab:tiendo}
    \renewcommand{\arraystretch}{1.3} % Tăng khoảng cách dòng cho dễ nhìn
    \begin{tabular}{|p{0.2\textwidth}|p{0.12\textwidth}|p{0.3\textwidth}|p{0.18\textwidth}|p{0.1\textwidth}|}
        \hline
        \textbf{Giai đoạn} & \textbf{Thời gian} & \textbf{Mô tả công việc cốt lõi} & \textbf{Kết quả đạt được} \\ \hline
        
        \textbf{Giai đoạn 1:} Khởi động \& Khảo sát & Tuần 1 & Thu thập tài liệu, khảo sát nghiệp vụ giao thông. Lựa chọn stack công nghệ và thiết kế kiến trúc. Setup Git và các công cụ làm việc. Họp thống nhất quy trình làm việc & Khởi tạo dự án thành công với kiến trúc Client-Server. \\ \hline
        
        \textbf{Giai đoạn 2:} Module Dashboard \& Thông tin Giao thông & Tuần 2 - Tuần 5 & Xử lý luồng ETL liên quan đến trạng thái giao thông hiện tại. Xử lý render 430.000 segment giao thông lên Mapbox map trên giao diện và các thông tin khác. & Luồng ETL Trafic Flow hoàn chỉnh cho TPHCM. Trang Dashboard Admin. \\ \hline
        
        \textbf{Giai đoạn 3:} Module Thống kê & Tuần 7 - Tuần 9 & Xây dựng các API liên quan đến thống kê. Phát triển UI \& chức năng Thống kê trên app dashboard. & Web Admin và App hoạt động ổn định. \\ \hline
        
        \textbf{Giai đoạn 4:} Module Dự báo \& Dẫn đường & Tuần 10 - Tuần 13 & Xây dựng \& train model dự báo kẹt xe trong 15 phút tới và tích hợp lên giao diện. Áp dụng thuật toán A* 2 chiều vào chức năng dẫn đường và phát triển giao diện mobile cho người dân. & Web Admin, Mobile App và hệ thống hoạt động ổn định. \\ \hline
        
        \textbf{Giai đoạn 5:} Triển khai \& Hoàn thiện báo cáo & Tuần 14 - Tuần 15 & Triển khai model, server lên Azure App Services. Triển khai giao diện lên Vercel. Viết tài liệu báo cáo đồ án, chuẩn bị slide và kịch bản Demo. & Quyển báo cáo. Slide thuyết trình. App production chạy ổn định. \\ \hline
    \end{tabular}
\end{table}
